% ════════════════════════════════════════════════════════════════════
% Chapter 7 — Mathematical Data Processing Pipeline
% ════════════════════════════════════════════════════════════════════
\section{Mathematical Data Processing Pipeline}\label{sec:pipeline}

The ``Golden Model'' pipeline is a multi-stage Python processing chain that transforms raw IMU telemetry into fully validated kinematic trajectories.
It is implemented as a sequence of numbered scripts (\file{step0}--\file{step8}) in the \file{golden\_model/} directory, each building upon the outputs of the previous stage.
This section presents the mathematical formulation of each stage.


\subsection{Step 0: Data Loading and Visualisation}

Raw CSV files are loaded and aligned.
The IMU time base is derived from the ESP32's monotonic microsecond counter (\code{esp\_timestamp\_us}) rather than the host timestamp, because the host timestamps arrive in bursts of 8 (one per ESP-NOW batch) separated by $\sim$\SI{8}{\milli\second} gaps, which would cause numerical instability in \code{np.gradient} and integration routines.
The ESP32 timestamp is anchored to the first IMU sample's host wall-clock time to share the camera's temporal reference.


\subsection{Step 1: Attitude Estimation and Gravity Removal}\label{sec:attitude}

The accelerometer measures the vector sum of gravity and the barbell's true linear acceleration in the \emph{sensor body frame}.
To isolate the linear acceleration, we must:
\begin{enumerate}[nosep]
  \item Continuously estimate the sensor's orientation (attitude) relative to the world frame.
  \item Rotate the measured acceleration into the world frame.
  \item Subtract the known gravity vector.
\end{enumerate}

Orientation is represented as a unit quaternion $\mathbf{q} = [q_w,\; q_x,\; q_y,\; q_z]^T$, avoiding the Gimbal lock singularity inherent to Euler-angle representations.

\subsubsection{Rotation via Quaternions}

The rotation matrix $\mathbf{R}(\mathbf{q})$ that transforms a vector from the body frame to the world frame is:
\begin{equation}\label{eq:rotmat}
  \mathbf{R}(\mathbf{q}) = \begin{bmatrix}
    1 - 2(q_y^2 + q_z^2) & 2(q_x q_y - q_w q_z) & 2(q_x q_z + q_w q_y) \\
    2(q_x q_y + q_w q_z) & 1 - 2(q_x^2 + q_z^2) & 2(q_y q_z - q_w q_x) \\
    2(q_x q_z - q_w q_y) & 2(q_y q_z + q_w q_x) & 1 - 2(q_x^2 + q_y^2)
  \end{bmatrix}
\end{equation}

The world-frame linear acceleration is then:
\begin{equation}\label{eq:grav_removal}
  \mathbf{a}_w = \mathbf{R}(\mathbf{q}) \, \mathbf{a}_b - \begin{bmatrix} 0 \\ 0 \\ g \end{bmatrix}
\end{equation}
where $\mathbf{a}_b$ is the raw accelerometer reading in the body frame and $g = \SI{9.80665}{\metre\per\second\squared}$.

\subsubsection{Initial Orientation from Gravity Vector}

The initial quaternion $\mathbf{q}_0$ is computed from the mean acceleration over the first 300 static samples using the TRIAD method:
\begin{align}
  \text{pitch} &= -\arcsin(\bar{a}_{x,0}) \\
  \text{roll}  &= \arcsin\!\left(\frac{\bar{a}_{y,0}}{\cos(\text{pitch})}\right) \\
  \text{yaw}   &= 0 \quad \text{(undefined without magnetometer)}
\end{align}

The Euler angles are converted to a quaternion via the ZYX convention.

\subsubsection{Madgwick AHRS Filter}

The Madgwick filter~\cite{madgwick2011} fuses gyroscope and accelerometer data using a gradient-descent optimisation that minimises the error between the predicted and measured gravity direction.
The quaternion rate of change is:
\begin{equation}\label{eq:madgwick}
  \dot{\mathbf{q}} = \frac{1}{2}\mathbf{q} \otimes \boldsymbol{\omega} - \beta \frac{\nabla f}{\|\nabla f\|}
\end{equation}
where $\boldsymbol{\omega}$ is the gyroscope angular velocity, $f$ is the objective function measuring the discrepancy between the measured and predicted gravity direction, and $\beta$ is the filter gain.

Our implementation uses an \textbf{adaptive $\beta$}: when $\left|\|\mathbf{a}\| - 1\right| < 0.1$ (i.e., the measured acceleration is close to $1\,g$, indicating the sensor is approximately stationary or in uniform motion), $\beta = 0.15$ for aggressive gravity correction; otherwise, $\beta = 0.02$ to avoid corrupting the attitude estimate during high-acceleration dynamic phases.

\subsubsection{Multiplicative Extended Kalman Filter (MEKF)}

As a cross-validation, the pipeline implements a 6-state MEKF with states $[\delta\boldsymbol{\theta}^T,\; \delta\mathbf{b}_g^T]^T$ representing the attitude error and gyroscope bias error.
The prediction step propagates the nominal quaternion via gyroscope integration:
\begin{equation}
  \mathbf{q}_{k+1} = \mathbf{q}_k \otimes \Delta\mathbf{q}(\boldsymbol{\omega} - \mathbf{b}_g, \Delta t)
\end{equation}
The update step observes the normalised accelerometer vector and corrects the attitude error via a multiplicative quaternion update, with the innovation gated by $\left|\|\mathbf{a}\| - 1\right| < 0.2$ to reject dynamic-acceleration phases.

\subsubsection{Versatile Quaternion-based Filter (VQF)}

The VQF~\cite{laidig2023} is the most sophisticated attitude estimator in the pipeline.
Unlike the Madgwick filter's single-gain structure, VQF implements a three-step architecture:
\begin{enumerate}[nosep]
  \item Gyroscope-only quaternion integration.
  \item Accelerometer-based inclination correction (independent of angular velocity).
  \item Optional magnetometer-based heading correction (not used in our application).
\end{enumerate}

VQF's key advantage for VBT is its separate \textbf{rest detection} module: it identifies periods of low motion using configurable thresholds on angular velocity and acceleration variance, enabling faster gyroscope bias convergence during the brief stationary phases between repetitions.


\subsection{Step 2: Velocity Integration}

The world-frame vertical acceleration $a_{w,z}$ is numerically integrated using the trapezoidal rule:
\begin{equation}
  v_z(t_{k+1}) = v_z(t_k) + a_{w,z}(t_k) \cdot \Delta t_k
\end{equation}

Without drift mitigation, the velocity diverges rapidly.
At this stage, simple linear drift removal is applied between manually identified rest boundaries:
\begin{equation}
  v_{\text{corr},k} = v_{\text{raw},k} - k \cdot \frac{v_{\text{raw},N}}{N}
\end{equation}
where $N$ is the segment length and $v_{\text{raw},N}$ is the raw velocity at the segment end.


\subsection{Step 3: Autonomous ZUPT Detection}

Zero-Velocity Updates (ZUPT) are the primary mechanism for bounding integration drift~\cite{skog2010}.
The autonomous ZUPT detector identifies stationary periods using a sliding-window variance test:
\begin{equation}
  \sigma^2_{a_z}(t) = \text{Var}\!\left[a_{w,z}(t-W/2 : t+W/2)\right]
\end{equation}
where $W = \SI{200}{\milli\second}$ (\SI{200}{} samples at \SI{1}{\kilo\hertz}).
A sample is classified as static if $\sigma^2_{a_z}(t) < \SI{0.05}{\metre\squared\per\second^4}$.
Isolated static detections shorter than \SI{300}{\milli\second} are suppressed as false positives.


\subsection{Step 7: Camera-Oracle ZUPT Strategies}\label{sec:oracle_zupt}

To isolate the bottleneck between ZUPT detection error and orientation drift, the pipeline implements a ``camera oracle'' experiment that bypasses the IMU stillness detector entirely.
Five strategies of ascending fusion are evaluated:

\begin{table}[H]
  \centering
  \caption{Camera-oracle ZUPT integration strategies.}
  \label{tab:zupt_strategies}
  \begin{tabular}{clp{7cm}}
    \toprule
    \textbf{ID} & \textbf{Name} & \textbf{Description} \\
    \midrule
    A & Velocity ZUPT only             & Force $v = 0$ at camera-detected rest anchors. Linear drift removal between anchors. \\
    B & Velocity + start-position      & As A, plus anchor the initial position to \code{cam\_pos\_i[anchor\_0]}. \\
    C & Full position + velocity anchor & Force $v = 0$ \emph{and} $p = p_\text{cam}$ at every rest anchor. \\
    D & Dense anchors (every \SI{200}{\milli\second}) & Pin $v = 0$ and $p = p_\text{cam}$ every \SI{200}{\milli\second}, regardless of motion. \\
    E & Rest-rest + \SI{50}{\milli\second} position chunks & Outer $v = 0$ boundaries at rest anchors; internal \SI{50}{\milli\second} position pins to camera. \\
    \bottomrule
  \end{tabular}
\end{table}

Camera-oracle rest anchors are located by identifying intervals where $|v_{z,\text{cam}}| < \SI{0.05}{\metre\per\second}$ for at least \SI{100}{\milli\second}.
The results of this experiment directly motivated the development of the 15-state ESKF.


\subsection{Step 8: 15-State Error-State Kalman Filter}\label{sec:eskf}

The ESKF is the culmination of the processing pipeline, achieving sub-\SI{50}{\milli\metre} position RMSE using only IMU data (no camera fusion) with lifting-specific ZUPT anchors.

\subsubsection{State Vector}

The nominal state comprises 16 storage elements:
\begin{equation}
  \mathbf{x} = [\underbrace{\mathbf{p}^T}_{3},\; \underbrace{\mathbf{v}^T}_{3},\; \underbrace{\mathbf{q}^T}_{4},\; \underbrace{\mathbf{b}_a^T}_{3},\; \underbrace{\mathbf{b}_g^T}_{3}]^T
\end{equation}
The error state $\delta\mathbf{x} \in \mathbb{R}^{15}$ uses the minimal rotation parameterisation $\delta\boldsymbol{\theta} \in \mathbb{R}^3$:
\begin{equation}
  \delta\mathbf{x} = [\delta\mathbf{p}^T,\; \delta\mathbf{v}^T,\; \delta\boldsymbol{\theta}^T,\; \delta\mathbf{b}_a^T,\; \delta\mathbf{b}_g^T]^T
\end{equation}

\subsubsection{Process Model}

The continuous-time dynamics are:
\begin{align}
  \dot{\mathbf{p}} &= \mathbf{v} \\
  \dot{\mathbf{v}} &= \mathbf{R}(\mathbf{q})(\mathbf{a}_m - \mathbf{b}_a) + \mathbf{g}_w \label{eq:vdot}\\
  \dot{\mathbf{q}} &= \tfrac{1}{2}\,\mathbf{q} \otimes (\boldsymbol{\omega}_m - \mathbf{b}_g) \\
  \dot{\mathbf{b}}_a &= \boldsymbol{\eta}_a \sim \mathcal{N}(\mathbf{0},\; \sigma_{ba}^2 \mathbf{I}_3) \\
  \dot{\mathbf{b}}_g &= \boldsymbol{\eta}_g \sim \mathcal{N}(\mathbf{0},\; \sigma_{bg}^2 \mathbf{I}_3)
\end{align}
where $\mathbf{g}_w = [0,\; 0,\; -g]^T$ is the gravity vector in the world frame and $\boldsymbol{\eta}$ represents random walk noise.

\subsubsection{Error-State Transition Matrix}

The linearised error-state transition matrix $\mathbf{F} = \mathbf{I}_{15} + \Delta t \cdot \mathbf{F}_c$ contains the following active blocks:
\begin{align}
  \mathbf{F}_{pv}   &= \mathbf{I}_3 \cdot \Delta t & &\text{(position driven by velocity)} \\
  \mathbf{F}_{v\theta} &= -[\mathbf{R}\,\mathbf{a}_c]_\times \cdot \Delta t & &\text{(velocity driven by attitude error)} \\
  \mathbf{F}_{v,ba}  &= -\mathbf{R} \cdot \Delta t & &\text{(velocity driven by accel bias)} \\
  \mathbf{F}_{\theta\theta} &= \mathbf{I}_3 - [\boldsymbol{\omega}_c]_\times \cdot \Delta t & &\text{(attitude error propagation)} \\
  \mathbf{F}_{\theta,bg} &= -\mathbf{I}_3 \cdot \Delta t & &\text{(attitude driven by gyro bias)}
\end{align}
where $[\cdot]_\times$ denotes the skew-symmetric matrix and $\mathbf{a}_c = \mathbf{a}_m - \mathbf{b}_a$.

\subsubsection{ZUPT Observation Model}

At rest-flagged samples, we observe $\mathbf{v} = \mathbf{0}$:
\begin{equation}
  \mathbf{y}_v = -\mathbf{v}_k, \quad \mathbf{H}_v = [\mathbf{0}_{3\times3},\; \mathbf{I}_3,\; \mathbf{0}_{3\times9}]
\end{equation}
The Kalman gain:
\begin{equation}
  \mathbf{K} = \mathbf{P}\,\mathbf{H}_v^T\,(\mathbf{H}_v\,\mathbf{P}\,\mathbf{H}_v^T + \mathbf{R}_v)^{-1}
\end{equation}
and the error-state correction $\delta\mathbf{x} = \mathbf{K}\,\mathbf{y}_v$ is injected into the nominal state:
\begin{align}
  \mathbf{p}_k &\leftarrow \mathbf{p}_k + \delta\mathbf{p} \\
  \mathbf{v}_k &\leftarrow \mathbf{v}_k + \delta\mathbf{v} \\
  \mathbf{q}_k &\leftarrow \mathbf{q}_k \otimes \delta\mathbf{q}(\delta\boldsymbol{\theta}) \\
  \mathbf{b}_a &\leftarrow \mathbf{b}_a + \delta\mathbf{b}_a \\
  \mathbf{b}_g &\leftarrow \mathbf{b}_g + \delta\mathbf{b}_g
\end{align}

The cross-covariance structure of $\mathbf{P}$ is the key mechanism: when velocity is forced to zero, the correction propagates through the $\delta\mathbf{v}$--$\delta\boldsymbol{\theta}$, $\delta\mathbf{v}$--$\delta\mathbf{b}_a$, and $\delta\mathbf{v}$--$\delta\mathbf{b}_g$ cross-terms, actively \emph{debiasing the sensors} as a byproduct of the velocity constraint.

\subsubsection{Zero Angular Rate Update (ZARU)}

During rest, we additionally observe that the true angular rate should be approximately zero:
\begin{equation}
  \mathbf{y}_g = \mathbf{b}_g - \boldsymbol{\omega}_m, \quad \mathbf{H}_g = [\mathbf{0}_{3\times12},\; -\mathbf{I}_3]
\end{equation}
This accelerates gyroscope bias convergence during the brief rest phases.

\subsubsection{Zero Acceleration Update (ZAU)}

At rest, the world-frame linear acceleration should be zero:
\begin{equation}
  \mathbf{a}_{\text{pred}} = \mathbf{R}\,(\mathbf{a}_m - \mathbf{b}_a) + \mathbf{g}_w \overset{!}{=} \mathbf{0}
\end{equation}
Linearising around the current attitude:
\begin{equation}
  \mathbf{H}_a = [-\mathbf{R}\,[\mathbf{a}_c]_\times,\; -\mathbf{R}] \in \mathbb{R}^{3\times6} \quad \text{(columns 6--14 of the 15-state)}
\end{equation}
This observation constrains the pitch and roll components of the attitude error, preventing gyroscope drift from misaligning the gravity subtraction vector.

\subsubsection{Lifting-Specific Near-Rest Detector}

Unlike pedestrian ZUPT, where foot-strike provides $>\SI{300}{\milli\second}$ of clean stationary data, a barbell only experiences ``rest'' for $\sim$\SI{50}{\milli\second} at the velocity zero-crossings at the top and bottom of each repetition.
VQF's default rest detector (\code{restMinT = 1.5\,s}) fires only $\sim$4 times in a 48-second squat set, starving the ESKF of updates.

Our custom detector accepts windows as short as \SI{50}{\milli\second} using two simultaneous conditions over a sliding window:
\begin{enumerate}[nosep]
  \item $\max_{w} |a_{w,z}| < \SI{0.5}{\metre\per\second\squared}$
  \item $\max_{w} |\boldsymbol{\omega}| < \SI{10}{\degree\per\second}$
\end{enumerate}
where $w$ is the window width.
The ESKF is designed to tolerate false-positive ZUPT detections gracefully: the observation noise $\mathbf{R}_v$ ensures that a spurious zero-velocity assertion during motion produces a large innovation that is appropriately down-weighted by the Kalman gain.

\subsubsection{Filter Noise Parameters}

\begin{table}[H]
  \centering
  \caption{ESKF noise parameters.}
  \label{tab:eskf_params}
  \begin{tabular}{llr}
    \toprule
    \textbf{Parameter} & \textbf{Symbol} & \textbf{Value} \\
    \midrule
    Accel measurement noise      & $\sigma_{a}$  & \SI{0.05}{\metre\per\second\squared} \\
    Gyro measurement noise       & $\sigma_{g}$  & \SI{0.005}{\radian\per\second} \\
    Accel bias random walk        & $\sigma_{ba}$ & \num{1e-4} \\
    Gyro bias random walk         & $\sigma_{bg}$ & \num{1e-5} \\
    ZUPT velocity observation     & $R_v$         & $(\num{0.001})^2$ \\
    ZARU gyro observation         & $R_g$         & $(\num{0.005})^2$ \\
    ZAU acceleration observation  & $R_{au}$      & $(\num{0.05})^2$ \\
    \midrule
    Initial $\delta\boldsymbol{\theta}$ uncertainty & $P_{\theta,0}$ & $(\SI{2}{\degree})^2$ \\
    Initial $\delta\mathbf{b}_a$ uncertainty        & $P_{ba,0}$     & $(\SI{0.05}{\metre\per\second\squared})^2$ \\
    Initial $\delta\mathbf{b}_g$ uncertainty        & $P_{bg,0}$     & $(\SI{0.5}{\degree\per\second})^2$ \\
    \bottomrule
  \end{tabular}
\end{table}


\subsection{Repetition Segmentation in the Golden Model}

The golden model pipeline segments repetitions by matching confirmed extrema (TOP/BOTTOM) in the ESKF-corrected position trace against the camera ground truth.
For each annotated repetition, the pipeline computes:
\begin{itemize}[nosep]
  \item Camera ROM vs.\ IMU ROM (from the max--min position within the rep boundaries).
  \item Per-rep ROM error in millimetres and percentage.
  \item Aggregate ROM RMSE and MAE across the set.
\end{itemize}
